\subsection{Security of RAG-Enabled Coding Agents}

\noindent \textbf{Poisoning of Retrieved Knowledge Sources.}
Recent work shows that retrieval-augmented generation (RAG) systems
can be attacked by poisoning the external content they retrieve.
\textsc{PoisonedRAG}~\cite{zou2025poisonedrag} studies knowledge corruption attacks
in which an adversary modifies the retrieval database so that the model produces
an attacker-chosen answer for a target query.
\textsc{AgentPoison}~\cite{chen2024agentpoison} further studies poisoning attacks
against agents' long-term memory or RAG knowledge bases across several application domains,
including autonomous driving, knowledge-intensive question answering, and healthcare.
Together, these works show that once an agent relies on external retrieved knowledge,
the knowledge source itself becomes a practical attack surface.

\noindent \textbf{Black-Box Attacks on Transparent RAG Pipelines.}
\textsc{RAG Paradox}~\cite{choi2025rag} identifies a structural
weakness in RAG systems that explicitly expose retrieved sources and content to users.
Although this transparency is intended to improve trust,
it also helps attackers observe which sources are used and how retrieved evidence is phrased,
making it easier to craft poisoned documents that are likely to be retrieved later.
This threat is directly relevant to coding agents,
because many practical agent frameworks reveal citations,
retrieved snippets, or tool outputs as part of their reasoning process.

\noindent \textbf{Defenses Against Retrieval Corruption.}
Recent defenses against RAG poisoning focus on consolidating knowledge
from multiple sources in order to reduce the influence of conflicting or malicious inputs.
\textsc{TrustRAG}~\cite{zhou2025trustrag} uses clustering to consolidate knowledge
from the model's internal state and external retrieved documents.
\textsc{Astute RAG}~\cite{wang2025astute} adaptively integrates internal and external knowledge,
using source-aware reasoning to resolve conflicts across different information sources.
\textsc{Certified Robust RAG}~\cite{xiang2024certifiably} proposes an isolate-then-aggregate strategy
that generates a separate response for each retrieved passage and then combines these
responses using a robust aggregation procedure.
However, these defenses are evaluated mainly in relatively simple tasks such as 
question answering.
It remains an open question how to design effective defenses for coding agents such as 
Codex~\cite{openai_codex_2026} or OpenCode~\cite{opencode_2026},
which typically require complex reasoning, 
and may execute code or take external actions.
These settings make compromise substantially more consequential 
and more difficult to defend against
than single-answer RAG tasks.
